À minha esposa, pelo companheirismo, pela paciência nos dias e noites dedicados a este trabalho e por acreditar em mim mesmo quando eu duvidava. Esta conquista também é sua.

Aos meus pais, pelo amor, pelos sacrifícios e pelo exemplo que me trouxeram até aqui. Tudo o que construí tem como base o que vocês me ensinaram.

Ao meu irmão, pela parceria de sempre, pelo incentivo e pelos momentos que tornaram a caminhada mais leve.

Ao Prof. Dr. Alan Michell Barros Alexandre, pela orientação, pela disponibilidade e pelas contribuições ao longo do desenvolvimento deste trabalho.

Aos membros da banca examinadora, pelo tempo dedicado à leitura do trabalho e pelas sugestões.

Aos meus colegas de curso, pela amizade, pelas horas de estudo compartilhadas e pela ajuda mútua nos momentos difíceis da graduação.

À Fundação Cearense de Meteorologia e Recursos Hídricos (FUNCEME) e à Companhia de Gestão dos Recursos Hídricos (COGERH), pela disponibilização das séries de vazões, das curvas cota-área-volume e dos planos de gestão proativa de secas utilizados neste trabalho.

Aos professores e servidores do Campus de Crateús da Universidade Federal do Ceará, pela formação recebida.
